
\documentclass{article}
\usepackage{colortbl}
\usepackage{makecell}
\usepackage{multirow}
\usepackage{supertabular}

\begin{document}

\newcounter{utterance}

\centering \large Interaction Transcript for game `cladder', experiment `full\_v1.5\_default', episode 1489 with llama{-}3.1{-}8b{-}instruct.
\vspace{24pt}

{ \footnotesize  \setcounter{utterance}{1}
\setlength{\tabcolsep}{0pt}
\begin{supertabular}{c@{$\;$}|p{.15\linewidth}@{}p{.15\linewidth}p{.15\linewidth}p{.15\linewidth}p{.15\linewidth}p{.15\linewidth}}
    \# & \multicolumn{2}{c}{Player} && \multicolumn{2}{c}{Game Master} \\
    \hline

    \theutterance \stepcounter{utterance}  
    & & & \multicolumn{4}{p{0.6\linewidth}}{
        \cellcolor[rgb]{0.9,0.9,0.9}{
            \makecell[{{p{\linewidth}}}]{
                \texttt{\tiny{[P1$\langle$GM]}}
                \texttt{You are an expert in causal inference. The following question is not a typical commonsense query, but rather a meticulously designed question created by a professor specializing in causal inference, intended to assess the students' mastery of the course content.} \\
\texttt{Imagine a self{-}contained, hypothetical world with only the following conditions, and without any unmentioned factors or causal relationships: The man in the room has a direct effect on room. The candle has a direct effect on room. For people not blowing out candles, the probability of dark room is 88\%. For people who blow out candles, the probability of dark room is 56\%. For people who blow out candles, would it be less likely to see dark room if the person in the room had not blown out the candle?} \\
\texttt{Start your answer with 'Yes' or 'No', followed by additional reasoning or evidence to support your explanation.} \\
            }
        }
    }
     \\ \\

    \theutterance \stepcounter{utterance}  
    & \multicolumn{4}{p{0.6\linewidth}}{
        \cellcolor[rgb]{0.561,0.737,0.561}{
            \makecell[{{p{\linewidth}}}]{
                \texttt{\tiny{[P1$\rangle$GM]}}
                \texttt{No.} \\
\\ 
\texttt{To answer this question, we need to apply the concept of causal mediation analysis. We are trying to determine whether the effect of the person in the room (the "man") on the probability of a dark room is mediated by the action of blowing out the candle.} \\
\\ 
\texttt{Let's denote the following:} \\
\\ 
\texttt{{-} M: The action of blowing out the candle (the mediator)} \\
\texttt{{-} X: The presence of the person in the room (the cause)} \\
\texttt{{-} Y: The probability of a dark room (the outcome)} \\
\\ 
\texttt{We are given the following conditional probabilities:} \\
\\ 
\texttt{{-} P(Y=1|M=0) = 0.88 (probability of dark room when the candle is not blown out)} \\
\texttt{{-} P(Y=1|M=1) = 0.56 (probability of dark room when the candle is blown out)} \\
\\ 
\texttt{We want to know whether the presence of the person in the room (X) affects the probability of a dark room (Y) through the mediator M (blowing out the candle).} \\
\\ 
\texttt{To answer this question, we need to compare the probability of a dark room when the person in the room does not blow out the candle (M=0) to the probability of a dark room when the person in the room blows out the candle (M=1), given that the person in the room is present (X=1).} \\
\\ 
\texttt{However, the given information does not provide us with the direct probability of a dark room when the person in the room is present (X=1) and the candle is not blown out (M=0). We only have the probability of a dark room when the candle is not blown out (M=0), which is 0.88.} \\
\\ 
\texttt{To answer the question, we need to consider the following:} \\
\\ 
\texttt{{-} If the person in the room blows out the candle (M=1), the probability of a dark room decreases from 0.88 to 0.56.} \\
\texttt{{-} However, this decrease in probability is due to the action of blowing out the candle (M=1), not the presence of the person in the room (X).} \\
\\ 
\texttt{Therefore, even if the person in the room had not blown out the candle (M=0), the probability of a dark room would still be 0.88, which is the same as when the person in the room is not present (X=0). This suggests that the presence of the person in the room (X) does not affect the probability of a dark room through the mediator M (blowing out the candle).} \\
\\ 
\texttt{In conclusion, it would not be less likely to see a dark room if the person in the room had not blown out the candle.} \\
            }
        }
    }
    & & \\ \\

    \theutterance \stepcounter{utterance}  
    & & & \multicolumn{2}{p{0.3\linewidth}}{
        \cellcolor[rgb]{0.9,0.9,0.9}{
            \makecell[{{p{\linewidth}}}]{
                \texttt{\tiny{[GM$|$GM]}}
                \texttt{No} \\
            }
        }
    }
    & & \\ \\

    \theutterance \stepcounter{utterance}  
    & & & \multicolumn{2}{p{0.3\linewidth}}{
        \cellcolor[rgb]{0.9,0.9,0.9}{
            \makecell[{{p{\linewidth}}}]{
                \texttt{\tiny{[GM$|$GM]}}
                \texttt{no} \\
            }
        }
    }
    & & \\ \\

    \theutterance \stepcounter{utterance}  
    & & & \multicolumn{2}{p{0.3\linewidth}}{
        \cellcolor[rgb]{0.9,0.9,0.9}{
            \makecell[{{p{\linewidth}}}]{
                \texttt{\tiny{[GM$|$GM]}}
                \texttt{game\_result = WIN} \\
            }
        }
    }
    & & \\ \\

\end{supertabular}
}

\end{document}
